% ===================================================================
%  12. TAULA — Geltokia. — Burdinbideak. — Ontziak.
%  Traduction 2026 d'apres texto/io, controlee sur texto/fr, texto/en
%  et texto/es. Consigne : voir texto/eu/10-taula-01.tex.
%
%  MEME OBSERVATION QU'AU TABLEAU 10 : le vocabulaire du rail est
%  arrive tout fait — « bagoia », « nasa », « semaforoa »,
%  « lokomotorra » — et ce qui reste basque est ce que le rail a trouve
%  en place : « maleta », « zamaketaria », « mahaizaina ».
% ===================================================================

%%K t12-apar-1 apar
\VUsaut{4.0mm}
\VUtitre{15.0pt}{60}{12. TAULA}
\VUsaut{2.4mm}
\VUpk{11.4pt}{40}{Geltokia. --- Burdinbideak. --- Ontziak.}
\VUsaut{3.4mm}
\textsc{Oharra.} --- \textit{Ordutegiak herrialde bakoitzean bereak dira, eta
horregatik eredutzat} \VUgras{taula frantsesaren} \textit{orduak hartu ditugu}.

%%K t12-01-1 p
1. --- Alemanian zeharreko bidaia bat egin berri dut. Abiatu aurretik
\VUgras{ordutegi-liburua} \textsuperscript{(15)} erosi nuen, trenen irteera-orduak
jakiteko. Trenik egokiena Paristik gaueko bederatzietan irteten zela eta
Estrasburgora ordu bata eta hamahiruetan iristen zela egiaztaturik, bidaia
prestatu nuen, eta hurrengo egunean geltokira eramanarazi nintzen.

%%K t12-02-1 p
2. --- Irteera-areto handian sartu nintzenean, herriko \VUgras{apaiz}
zahar baten \textsuperscript{(2)} atzetik, zeinak bere \VUgras{bidaia-poltsa}
\textsuperscript{(3)} eskuan baitzeraman, \VUgras{bidaiari} asko
\textsuperscript{(1)} bildu zen jada.

%%K t12-02-2 p
\VUgras{Telegrama} bat \textsuperscript{(5)} bidali behar nuenez,
\VUgras{kontrolatzaile} bati \textsuperscript{(6)} eskatu nion
\VUgras{telegrafo-bulegoa} \textsuperscript{(4)} erakusteko.

%%K t12-03-1 p
3. --- \VUgras{Itxarongelan} \textsuperscript{(7)} sartu aurretik, joan-etorriko
\VUgras{txartel} bat \textsuperscript{(9)} hartzera joan nintzen.

%%K t12-04-1 p
4. --- \VUgras{Leihatilan} \textsuperscript{(8)} ez nuen luze itxaron, jende
gutxi baitzegoen. \VUgras{Soldadu} batek \textsuperscript{(10)}, oporretara
zihoanak, adeitasunez utzi zidan bere txanda, eta hala astia izan nuen
\VUgras{geltokiburuari} \textsuperscript{(13)} galdera batzuk egiteko, zeina
zaintzako \VUgras{komisarioarekin} \textsuperscript{(12)} eta txandako
\VUgras{jendarmearekin} \textsuperscript{(11)} solasean baitzegoen.

%%K t12-tit-1 sub
\VUpk{10.2pt}{40}{Ekipajea ematea.}
\VUsaut{3.0mm}

%%K t12-05-1 p
5. --- \VUgras{Liburu-dendan} \textsuperscript{(14)} egunkari bat erosi ondoren,
nire \VUgras{ekipajea} \textsuperscript{(17)} pisatzeko agindu nuen ---
\VUgras{zamaketari} batek \textsuperscript{(16)} ekarri berri zuen
\VUgras{ekipaje-organ} \textsuperscript{(19)}. \VUgras{Funtzionarioak}
\textsuperscript{(22)}, \VUgras{balantzaren} \textsuperscript{(21)} ardura
duenak, esan zidan nire kutxak hogeita hamabost kilo pisatzen
zituela; horrek bost kiloko \VUgras{gainpisua} egiten zuen, zeinagatik gehigarria
ordaindu behar izan bainuen ekipajeen \VUgras{leihatilan}
\textsuperscript{(23)}.

%%K t12-06-1 p
6. --- Goizegi iritsi nintzenez, astia izan nuen geltokiko bidaiarien mugimenduari
begiratzeko. Batzuek beren esku-ekipajeak uzten edo jasotzen zituzten
\VUgras{gordailuan} \textsuperscript{(25)}; besteek \VUgras{ordutegiak}
\textsuperscript{(26)} kontsultatzen zituzten. Iristen ziren bidaiariak
\VUgras{irteerara} \textsuperscript{(32)} presatzen ziren beren \VUgras{maleta}
\textsuperscript{(24)} eskuan. Batzuk \VUgras{ekipaje-jasotzean}
\textsuperscript{(27)} zain zeuden eta beren \VUgras{txartela}
\textsuperscript{(29)} erakusten zuten beren kutxak, \VUgras{kaxak}
\textsuperscript{(28)} edo \VUgras{saskiak} \textsuperscript{(30)}
berreskuratzeko.

%%K t12-07-1 p
7. --- \VUgras{Zerga-funtzionarioek} \textsuperscript{(33)} arretaz aztertzen
zituzten fardelak, aitortzeko zerbait ote zegoen jakiteko.

%%K t12-08-1 p
8. --- Bidaiari batzuk \VUgras{jatetxera} \textsuperscript{(34)} zihoazen, non
\VUgras{mahaizain} batek \textsuperscript{(35)} arretaz zerbitzatzen baitzien.
Presaka ibili behar zuten, \VUgras{trena} \textsuperscript{(36)} korrikoa
baitzen eta ez baitzen luze geldituko geltokian.

%%K t12-09-1 p
9. --- Gero irteerako nasara atera nintzen eta \VUgras{bagoi} zuzen bat
\textsuperscript{(37)} bilatu nuen, bidean aldatu behar ez izateko. Korridoredun
bagoi batean sartu nintzen, non erretzaileentzako \VUgras{konpartimentu} bat
\textsuperscript{(38)} baitzegoen eta beste bat «bakarrik doazen andreentzat»
gordea.

%%K t12-tit-2 sub
\VUpk{10.2pt}{40}{Gaueko irteera.}
\VUsaut{3.0mm}

%%K t12-10-1 p
10. --- \VUgras{Estriboan} \textsuperscript{(40)} igo, \VUgras{atetxoa}
\textsuperscript{(39)} itxi, \VUgras{gortina} \textsuperscript{(41)} jaso eta
nire esku-ekipajeak \VUgras{saretan} \textsuperscript{(43)} jarri ondoren,
\VUgras{eserlekuan} \textsuperscript{(42)} eseri nintzen. \VUgras{Lanparazainak}
\textsuperscript{(44)} lanparak pizten ikusi nituenean, gogoratu nintzen gaua
bagoian igaro behar nuela, eta \VUgras{buruko} \textsuperscript{(45)}
alokatzeaz arduratzen den langilea deitu nuen, bat eskatzeko.

%%K t12-11-1 p
11. --- Trenaren kondaktorea etorri zen txartelak egiaztatzera. Galdetu nion
zergatik ez ginen oraindik abiatu, ordua iritsi baitzen. Erantzun zidan
\VUgras{tren-buruzagia} \textsuperscript{(46)} bizikletak \VUgras{ekipaje-bagoian}
\textsuperscript{(47)} sartzen ari zela. Gainera, ez zituzten oraindik
\VUgras{oin-berogailu} guztiak \textsuperscript{(48)} aldatu.

%%K t12-12-1 p
12. --- Baina laster abiatuko ginen, \VUgras{sugilea}
\textsuperscript{(50)} bere \VUgras{tenderrean} \textsuperscript{(49)} ikatza
hartzen ari baitzen \VUgras{lokomotorraren} \textsuperscript{(54)} sua
bizitzeko, eta \VUgras{makinistak} \textsuperscript{(51)}
\VUgras{geltokiburu-laguntzailearen} \textsuperscript{(52)} keinuaren zain
baitzegoen soilik, zeina \VUgras{nasan} \textsuperscript{(53)} baitzegoen.

%%K t12-13-1 p
13. --- Lokomotorraren \VUgras{galdarak} \textsuperscript{(56)} burrunba ilun bat
egiten zuen; \VUgras{tximiniak} \textsuperscript{(55)} ke-uholdeak botatzen
zituen \VUgras{lurrunarekin} \textsuperscript{(60)} nahasita. \VUgras{Txistu}
zorrotz bat \textsuperscript{(59)} entzun zen eta \VUgras{pistoiak}
\textsuperscript{(58)} mugitzen hasi ziren, makina astunaren \VUgras{gurpilak}
\textsuperscript{(57)} eraginez.

%%K t12-tit-3 sub
\VUsaut{3.0mm}
\VUpk{10.2pt}{40}{Goazen!}
\VUsaut{3.0mm}

%%K t12-14-1 p
14. --- Trena, \VUgras{errailen} \textsuperscript{(64)} gainean zihoala, gero eta
azkarrago zihoan; \VUgras{nasak} \textsuperscript{(65)} amaitu ziren;
\VUgras{komunak} \textsuperscript{(66)} igaro genituen, eta baita beste
\VUgras{bidean} \textsuperscript{(63)} geldirik zeuden bagoi lerro bat ere:
\VUgras{lo-bagoiak} \textsuperscript{(67)}, \VUgras{jatetxe-bagoi} bat
\textsuperscript{(68)}, \VUgras{posta-bagoi} bat \textsuperscript{(69)}.

%%K t12-noto-1 noto
\VUnotes{143.2mm}{%
\textit{(*) Oharra.} --- \textit{Bidaia, jakina, gerra aurreko mugan igarotzen zen
bezala deskribatu da.}%
}

%%K t12-15-1 p
15. --- Gero \VUgras{salgai-geltokia} \textsuperscript{(71)} igaro zen gure
begien aurretik. Teilapeetan langileek \VUgras{salgaiak}
\textsuperscript{(72)} kargatzen zituzten, abiadura txikian bidaltzen direnak.
\VUgras{Maniobra-langileek}
\textsuperscript{(76)}, \VUgras{ur-dorrearen} \textsuperscript{(73)}
ondokoek, \VUgras{abere-bagoiak} \textsuperscript{(75)} eta
\VUgras{plataformak} \textsuperscript{(79)} mugitzen zituzten
\VUgras{salgai-tren} bat \textsuperscript{(80)} osatzeko.

%%K t12-16-1 p
16. --- Azken \VUgras{orratza} \textsuperscript{(78)} igaro genuen, zeinaren
ondoan orratzaina baitzegoen; \VUgras{semaforoa} \textsuperscript{(86)} igaro
genuen eta geltokia urrunean desagertu zen.

%%K t12-17-1 p
17. --- Laster \VUgras{burdinazko zubira} \textsuperscript{(74)} atera ginen,
zeinak \VUgras{ibaia} \textsuperscript{(81)} zeharkatzen baitu. Zubi hori igaro
ondoren, ibaiaren ondotik jarraitu genuen, non \VUgras{atoiontzi} indartsuek
\textsuperscript{(82)} \VUgras{gabarra} astunak \textsuperscript{(83)}
tiratzen baitzituzten. \VUgras{Bidaiari-ontzi} bat ere \textsuperscript{(84)}
ikusi genuen, \VUgras{kai-mahairatzera} \textsuperscript{(85)} hurbiltzen ari
zena.

%%K t12-18-1 p
18. --- Geltoki batzuk abiadura izugarrian igaro ondoren, \VUgras{tunel}
batzuk \textsuperscript{(87)} zeharkatu genituen eta \VUgras{bide-zaindarien}
kontaezinezko \VUgras{etxetxoak} \textsuperscript{(88)} utzi genituen atzean.
Trenaren kulunkak lokartuta, lo sakonean erori nintzen.

%%K t12-tit-4 sub
\VUsaut{3.0mm}
\VUpk{10.2pt}{40}{Deutsch-Avricourt geltokia.}
\VUsaut{3.0mm}

%%K t12-19-1 p
19. --- Bat-batean langile baten ahotsak esnatu ninduen, oihuka:
«Deutsch-Avricourt! \textsuperscript{(*)} Denak jaitsi!» Aduanaren azterketa eta
mugako gogaikarrikeria guztiak izan ziren. Nire poltsikoko erlojua geltokiko
\VUgras{erloju handiaren} \textsuperscript{(89)} arabera jarri behar izan nuen,
zeina berrogeita hamabost minutu aurreratuta baitzihoan; ozta-ozta esnatuta,
nekez bereizten nuen \VUgras{orratz handia} \textsuperscript{(90)}
\VUgras{txikitik} \textsuperscript{(91)}; \VUgras{esfera} bera
\textsuperscript{(92)} erabat lausotuta zegoen goizeko lainoaz.

%%K t12-20-1 p
20. --- Aduanazainek beren azterketa egiten zuten bitartean, aharrausika
begiratzen nituen geltokiko hormak apaintzen dituzten \VUgras{kartelak}
\textsuperscript{(93)}. Gosaldu nuen denbora igarotzeko, Alemaniako
\VUgras{sarearen} \textsuperscript{(94)} mapa kontsultatu nuen Estrasburgotik
zenbat falta zitzaidan ikusteko, eta nire bagoira itzuli nintzen, laster helmugara
iristeko itxaropenak adoretua.
\VUsaut{6.0mm}
\VUfilet{22mm}
